% GENERATED FILE. Edit canonical lessons, then run npm run generate:books.
% canonical-source-hash: fnv1a64:bdb85550e67fd32b
% canonical-lessons: MR-C291-nisatne, MR-C291-sutka, MR-C291-sukhrup, MR-C291-bepatta, MR-C291-dhundalne

\chapter{A Lucky Escape}
\label{ch:mr-a-lucky-escape}
\hlchaptermodality{291}

\begin{chapteropening}
\textbf{By the end of this chapter:} \emph{I can say that someone escaped, was rescued, or went missing.}

Complete the last lesson of chapter 291: I can say that someone escaped, was rescued, or went missing.
\end{chapteropening}

\section[\texorpdfstring{\mr{निसटणे} (nisaṭṇe)}{निसटणे (nisaṭṇe)}]{\mr{निसटणे} (nisaṭṇe) — to slip away; to escape}

\label{lesson:MR-C291-nisatne}



\begin{quote}
Before the new one: say the Marathi for a prison, then the Marathi for a prisoner.
\end{quote}

\subsection*{You'll want to know: \mr{निसटणे}}
\textbf{\mr{निसटणे}} — \emph{nisaṭṇe} — \textquotedblleft{}to slip away; to escape\textquotedblright{}.

\begin{quote}
\mr{दार} \mr{उघडं} \mr{असतं} \mr{तर} \mr{मांजर} \mr{बाहेर} \mr{निसटलं} \mr{असतं}.
\end{quote}

\begin{quote}
\emph{dār ughḍaṁ astaṁ tar māñjar bāher nisaṭlaṁ astaṁ.}
\end{quote}

\begin{quote}
If the door had been open, the cat would have slipped out.
\end{quote}

Past (he …): \mr{तो} \mr{निसटला}, \emph{to nisaṭlā}.

The place left takes -\mr{तून}: \mr{खोलीतून} \mr{निसटणे}, to slip out of a room.

\begin{cousinweb}
\textbf{\mr{निसटणे}} (\emph{nisaṭṇe}) is inherited: CDIAL lists Marathi \emph{nisaṭṇe}, 'to slip away', under \emph{*niḥsṛṣṭa}, 'sent forth' (CDIAL 7132).
\end{cousinweb}

\begin{grammarlens}[title={What you've built}]
One more word for telling what happened.
\end{grammarlens}

\subsection*{Your turn}
\begin{itemize}
\raggedright
  \item \emph{Say it:} \emph{nisaṭṇe}
  \item \emph{Say it:} \emph{nisaṭṇe}, once more
  \item \emph{Say it:} \emph{dār ughḍaṁ astaṁ tar māñjar bāher nisaṭlaṁ astaṁ.} — If the door had been open, the cat would have slipped out.
  \item \emph{Recall:} say the Marathi for a prison, then the Marathi for a prisoner, then say \emph{nisaṭṇe} again
\end{itemize}

\begin{tcolorbox}[breakable,colback=teal!4,colframe=teal!35!black,title={Before you move on}]
What does \mr{निसटणे} mean? (\textquotedblleft{}To slip away; to escape\textquotedblright{}.) Say it once more.
\end{tcolorbox}
\section[\texorpdfstring{\mr{सुटका} (suṭkā)}{सुटका (suṭkā)}]{\mr{सुटका} (suṭkā) — a rescue; a release}

\label{lesson:MR-C291-sutka}



\begin{quote}
Before the new one: say the Marathi for a prisoner, then the Marathi for to slip away.
\end{quote}

\subsection*{You'll want to know: \mr{सुटका}}
\textbf{\mr{सुटका}} — \emph{suṭkā} — \textquotedblleft{}a rescue; a release\textquotedblright{}.

\begin{quote}
\mr{तीन} \mr{तासांनंतर} \mr{मुलाची} \mr{सुटका} \mr{झाली}.
\end{quote}

\begin{quote}
\emph{tīn tāsānnantar mulācī suṭkā jhālī.}
\end{quote}

\begin{quote}
After three hours the child was rescued.
\end{quote}

Feminine.

\mr{सुटका} \mr{होणे}, to be rescued or set free; \mr{सुटका} \mr{करणे}, to rescue.

\begin{cousinweb}
\textbf{\mr{सुटका}} (\emph{suṭkā}) is built on the stem \emph{suṭ-}, 'to come loose, be set free': CDIAL lists Marathi \emph{suṭṇẽ}, 'to get loose', and \emph{suṭī}, 'remission', under \emph{*kṣuṭyate}, 'is released' (CDIAL 3707).
\end{cousinweb}

\begin{grammarlens}[title={What you've built}]
One more word for telling what happened.
\end{grammarlens}

\subsection*{Your turn}
\begin{itemize}
\raggedright
  \item \emph{Say it:} \emph{suṭkā}
  \item \emph{Say it:} \emph{suṭkā}, once more
  \item \emph{Say it:} \emph{tīn tāsānnantar mulācī suṭkā jhālī.} — After three hours the child was rescued.
  \item \emph{Recall:} say the Marathi for a prisoner, then the Marathi for to slip away, then say \emph{suṭkā} again
\end{itemize}

\begin{tcolorbox}[breakable,colback=teal!4,colframe=teal!35!black,title={Before you move on}]
What does \mr{सुटका} mean? (\textquotedblleft{}A rescue; a release\textquotedblright{}.) Say it once more.
\end{tcolorbox}
\section[\texorpdfstring{\mr{सुखरूप} (sukhrūp)}{सुखरूप (sukhrūp)}]{\mr{सुखरूप} (sukhrūp) — safe and sound}

\label{lesson:MR-C291-sukhrup}



\begin{quote}
Before the new one: say the Marathi for to slip away, then the Marathi for a rescue.
\end{quote}

\subsection*{You'll want to know: \mr{सुखरूप}}
\textbf{\mr{सुखरूप}} — \emph{sukhrūp} — \textquotedblleft{}safe and sound\textquotedblright{}.

\begin{quote}
\mr{सगळे} \mr{प्रवासी} \mr{सुखरूप} \mr{घरी} \mr{पोहोचले}.
\end{quote}

\begin{quote}
\emph{sagḷe pravāsī sukhrūp gharī pohocle.}
\end{quote}

\begin{quote}
All the passengers reached home safe and sound.
\end{quote}

\mr{सुखरूप} keeps one form: \mr{तो} \mr{सुखरूप} \mr{आहे}, \mr{ती} \mr{सुखरूप} \mr{आहे}.

\begin{cousinweb}
\textbf{\mr{सुखरूप}} (\emph{sukhrūp}) is built from Sanskrit \emph{sukha}, 'ease, well-being', and \emph{rūpa}, 'form' (Monier-Williams): in a good state.
\end{cousinweb}

\begin{grammarlens}[title={What you've built}]
One more word for telling what happened.
\end{grammarlens}

\subsection*{Your turn}
\begin{itemize}
\raggedright
  \item \emph{Say it:} \emph{sukhrūp}
  \item \emph{Say it:} \emph{sukhrūp}, once more
  \item \emph{Say it:} \emph{sagḷe pravāsī sukhrūp gharī pohocle.} — All the passengers reached home safe and sound.
  \item \emph{Recall:} say the Marathi for to slip away, then the Marathi for a rescue, then say \emph{sukhrūp} again
\end{itemize}

\begin{tcolorbox}[breakable,colback=teal!4,colframe=teal!35!black,title={Before you move on}]
What does \mr{सुखरूप} mean? (\textquotedblleft{}Safe and sound\textquotedblright{}.) Say it once more.
\end{tcolorbox}
\section[\texorpdfstring{\mr{बेपत्ता} (bepattā)}{बेपत्ता (bepattā)}]{\mr{बेपत्ता} (bepattā) — missing}

\label{lesson:MR-C291-bepatta}



\begin{quote}
Before the new one: say the Marathi for a rescue, then the Marathi for safe and sound.
\end{quote}

\subsection*{You'll want to know: \mr{बेपत्ता}}
\textbf{\mr{बेपत्ता}} — \emph{bepattā} — \textquotedblleft{}missing\textquotedblright{}.

\begin{quote}
\mr{तो} \mr{मुलगा} \mr{दोन} \mr{दिवसांपासून} \mr{बेपत्ता} \mr{आहे}.
\end{quote}

\begin{quote}
\emph{to mulgā don divsāmpāsūn bepattā āhe.}
\end{quote}

\begin{quote}
That boy has been missing for two days.
\end{quote}

\mr{बेपत्ता} keeps one form: \mr{बेपत्ता} \mr{मुलगी}, \mr{बेपत्ता} \mr{माणूस}. \mr{पत्ता} is an address: with no address, no trace.

\begin{cousinweb}
\textbf{\mr{बेपत्ता}} (\emph{bepattā}) is built from the Persian prefix \emph{be-}, 'without', and \emph{pattā}, 'an address, a trace'; the origin of \emph{pattā} itself is not settled.
\end{cousinweb}

\begin{grammarlens}[title={What you've built}]
One more word for telling what happened.
\end{grammarlens}

\subsection*{Your turn}
\begin{itemize}
\raggedright
  \item \emph{Say it:} \emph{bepattā}
  \item \emph{Say it:} \emph{bepattā}, once more
  \item \emph{Say it:} \emph{to mulgā don divsāmpāsūn bepattā āhe.} — That boy has been missing for two days.
  \item \emph{Recall:} say the Marathi for a rescue, then the Marathi for safe and sound, then say \emph{bepattā} again
\end{itemize}

\begin{tcolorbox}[breakable,colback=teal!4,colframe=teal!35!black,title={Before you move on}]
What does \mr{बेपत्ता} mean? (\textquotedblleft{}Missing\textquotedblright{}.) Say it once more.
\end{tcolorbox}
\section[\texorpdfstring{\mr{धुंडाळणे} (dhuṇḍāḷṇe)}{धुंडाळणे (dhuṇḍāḷṇe)}]{\mr{धुंडाळणे} (dhuṇḍāḷṇe) — to comb through; to search high and low}

\label{lesson:MR-C291-dhundalne}



\begin{quote}
Before the new one: say the Marathi for safe and sound, then the Marathi for missing.
\end{quote}

\subsection*{You'll want to know: \mr{धुंडाळणे}}
\textbf{\mr{धुंडाळणे}} — \emph{dhuṇḍāḷṇe} — \textquotedblleft{}to comb through; to search high and low\textquotedblright{}.

\begin{quote}
\mr{आम्ही} \mr{हरवलेल्या} \mr{मुलासाठी} \mr{सगळं} \mr{गाव} \mr{धुंडाळलं}.
\end{quote}

\begin{quote}
\emph{āmhī haravlelyā mulāsāṭhī sagḷaṁ gāv dhuṇḍāḷlaṁ.}
\end{quote}

\begin{quote}
We searched the whole village for the lost child.
\end{quote}

Past, with -\mr{ने} (he …): \mr{त्याने} \mr{धुंडाळलं}, \emph{tyāne dhuṇḍāḷlaṁ}.

The place searched is the object: \mr{घर} \mr{धुंडाळणे}, to search the whole house.

\begin{cousinweb}
\textbf{\mr{धुंडाळणे}} (\emph{dhuṇḍāḷṇe}) is inherited: CDIAL lists Marathi \emph{dhuṇḍāḷṇe}, 'to search', under \emph{*dhūṇḍh}, 'seek' (CDIAL 6839).
\end{cousinweb}

\begin{grammarlens}[title={What you've built}]
One more word for telling what happened.
\end{grammarlens}

\subsection*{Your turn}
\begin{itemize}
\raggedright
  \item \emph{Say it:} \emph{dhuṇḍāḷṇe}
  \item \emph{Say it:} \emph{dhuṇḍāḷṇe}, once more
  \item \emph{Say it:} \emph{āmhī haravlelyā mulāsāṭhī sagḷaṁ gāv dhuṇḍāḷlaṁ.} — We searched the whole village for the lost child.
  \item \emph{Recall:} say the Marathi for safe and sound, then the Marathi for missing, then say \emph{dhuṇḍāḷṇe} again
\end{itemize}

\begin{tcolorbox}[breakable,colback=teal!4,colframe=teal!35!black,title={Before you move on}]
What does \mr{धुंडाळणे} mean? (\textquotedblleft{}To comb through; to search high and low\textquotedblright{}.) Say it once more.
\end{tcolorbox}
